\documentclass[10pt,a4paper]{amsart}
\usepackage[margin=19mm]{geometry}
\usepackage{amsmath,amssymb,mathtools}
\usepackage[scr=rsfs,cal=euler]{mathalfa}
\usepackage{xfrac}
\usepackage[unicode,hidelinks,bookmarks=false]{hyperref}
\newcommand{\ie}{i.e.\ }
\newcommand{\cf}{cf.\ }
\newcommand{\resp}{respectively\ }
\newcommand{\Subsection}{\subsection}
\providecommand{\nobreakdash}{\nobreak\hbox{-}}
\setlength{\emergencystretch}{2em}
\multlinegap0pt
\usepackage{tikz}
\usepackage{float}
\usepackage[normalem]{ulem}
\usepackage{amssymb}
\newtheorem{lemma}{Lemma}
\def\B{\mathcal{B}}
\def \IR {\mathbb{R}}
\def \Th {{\mathcal T}_h}
\def \d {{\rm d}}
\def\dx{\,\d x }
\def\gradn{\nabla_{\!\!\perp}}
\def\gradt{\nabla_{\!\!{\scriptscriptstyle\parallel}}}
\def\tD{\widetilde{D}}
\def\tu{\tilde{u}}
\def\width{\bar{h}}
\title{The tempered finite element method}
\author{Antoine {Quiriny}, Václav {Kučera}, Jonathan {Lambrechts}, Nicolas {Moës}, Jean-François {Remacle}}
\date{Source: arXiv:2411.17564v2 (CC BY 4.0)}
\begin{document}
\maketitle

\noindent\emph{Source and licence.} The tempered finite element method. Version arXiv:2411.17564v2 (CC BY 4.0), distributed by arXiv under the Creative Commons Attribution 4.0 International licence, http://creativecommons.org/licenses/by/4.0/, as recorded in the arXiv licence metadata for this exact version and shown on the abstract page https://arxiv.org/abs/2411.17564v2. Full licence notice: \texttt{licenses/arxiv-2411.17564-CC-BY-4.0.txt}.\par\smallskip

\noindent\textit{Editorial scope.} Counted unit: the article's lemma lem:FEM\_zero (source0.tex lines 2559--2571). The article's complete original proof of it is delivered with it (source0.tex lines 2588--2624, one \texttt{proof} environment of 37 lines, which the article places away from the statement). No mathematics is written, repaired or re-derived: the statement and the proof are the article's own lines, in the article's own notation, and no retained line is substituted or re-flowed.  Retained uncounted as the source-local context the proof needs: lines 558--561; lines 2553--2556.  Nothing was omitted from any retained span, so the package carries no omission record.  No retained line contains a citation, so the wrapper carries no bibliography and no bibliography entry.  No retained line contains a figure environment, an image-inclusion command or drawing code, so no image is delivered and the package declares no asset dependency.  Editorial formatting only, in addition: an AMS \texttt{amsart} preamble that restates the article's own declarations and macros used by the retained lines, namely the environments \texttt{lemma} on separate counters, together with the standard packages those lines need.  The retained environments print in this document as Lemma 1 in order of appearance, whereas the article prints the same items with the numbering of its own full document.  The complete native source of the article as distributed in the arXiv e-print for this exact version is bundled unchanged as \texttt{sources/169163/source0.tex}.
\begin{equation}
\int_\Omega d(x)\nabla \tu_h\cdotp\nabla v_h \dx=(f,v_h),\quad\forall v_h\in V_h^0.
\label{eq:FEM_mod}
\end{equation}

\begin{equation}
\label{eq:FEM_zero:3}
\int_{\Omega\setminus\B} \nabla \tu_h\cdotp\nabla v_h \dx +D\int_{\B} \gradn \tu_h \cdotp\gradn v_h +\gradt \tu_h\cdotp\gradt v_h \dx=(f,v_h).
\end{equation}

\begin{lemma}
\label{lem:FEM_zero}
Let $D=\width\tD$. Then the tempered FEM scheme \eqref{eq:FEM_mod} for $\width=0$ is equivalent to the formulation
\begin{equation}
\label{eq:FEM_zero}
\int_{\Omega\setminus\Gamma} \nabla \tu_h\cdotp\nabla v_h \dx +\tD\frac{h}{2}\sum_{x_i\in\Gamma}[\tu_h]_{x_i}[v_h]_{x_i}=(f,v_h),\quad\forall v_h\in V_h^\Gamma,
\end{equation}
where
\begin{equation*}
V_h^\Gamma = \{v_h\in C(\Omega\setminus\Gamma):\ v_h|_K\in P^1(K)\,\forall K\in\Th, v_h|_{\partial\Omega}=0\}
\end{equation*}
is the set of globally continuous piecewise linear functions discontinuous on $\Gamma$.
\end{lemma}

\begin{proof}[Proof of Lemma \ref{lem:FEM_zero}]
Consider the scheme \eqref{eq:FEM_mod} in the form \eqref{eq:FEM_zero:3}. Let $\varphi_i$, $i=1,\ldots,M$ be the standard `tent' basis functions of $V_h$ corresponding to individual vertices of $\Th$. We write $\tu_h(x)=\sum_{j=1}^M U_j\varphi_j(x)$, where $U_j\in \IR$, and set $v_h:=\varphi_k$. Then, \eqref{eq:FEM_zero:3} becomes
\begin{equation}
\label{lem:FEM_zero:1}
\sum_{j=1}^M U_j\int_{\Omega\setminus\B} \nabla \varphi_j\cdotp\nabla \varphi_k \dx +\sum_{j=1}^M U_j\width\tD\int_{\B} \gradn \varphi_j \cdotp\gradn \varphi_k +\gradt \varphi_j\cdotp\gradt \varphi_k \dx=(f,\varphi_k).
\end{equation}
Now consider that for all $j$, it holds $|\gradt \varphi_j|\leq1/h$ on each $K\in\B$ and that $|\B|\leq L\width$. Therefore, for fixed $h$,
\begin{equation*}
\label{lem:FEM_zero:2}
\Big|\width\tD\int_{\B} \gradt \varphi_j\cdotp\gradt \varphi_k \dx\Big| \leq \width^2 L\tD h^{-2} \to 0,\text{ as }\width\to 0^+,
\end{equation*}
therefore these terms `disappear' from the limiting system of linear equation (\ref{lem:FEM_zero:1}) for $\width=0$.

On the other hand, on an element $K_i\in\B$ with apex $x_i$ we have
\begin{equation*}
\gradn \varphi_j=\pm\frac{\varphi_j(x_i)-\varphi_j(p_i)}{\width}n_\Gamma,
\end{equation*}
for all $j$ (the sign depends on the orientation of $K$ with respect to $n_\Gamma$). Therefore (since $|K|=\tfrac{1}{2}\width h$ for all $K\in\B$),
\begin{equation}
\nonumber
\label{lem:FEM_zero:3}
\begin{split}
&\width\tD\int_{\B} \gradn \varphi_j \cdotp\gradn \varphi_k \dx= \width\tD\sum_{K_i\in\B}\frac{1}{2}\width h \frac{\varphi_j(x_i)-\varphi_j(p_i)}{\width}
\frac{\varphi_k(x_i)-\varphi_k(p_i)}{\width}\\
&\ = \tD\frac{h}{2}\sum_{K_i\in\B} \big(\varphi_j(x_i)-\varphi_j(p_i)\big) \big(\varphi_k(x_i)-\varphi_k(p_i)\big)\longrightarrow \tD\frac{h}{2}\sum_{x_i\in\Gamma} [\varphi_j]_{x_i}[\varphi_k]_{x_i},
\end{split}
\end{equation}
for $\width\to 0^+$.

Altogether, the limiting formulation of (\ref{lem:FEM_zero:1}) for $\width=0$ is 
\begin{equation*}
\label{lem:FEM_zero:4}
\sum_{j=1}^M U_j\int_{\Omega\setminus\Gamma} \nabla \varphi_j\cdotp\nabla \varphi_k \dx +\sum_{j=1}^M U_j\tD\frac{h}{2} \sum_{x_i\in\Gamma} [\varphi_j]_{x_i}[\varphi_k]_{x_i} =(f,\varphi_k),
\end{equation*}
which is equivalent to \eqref{eq:FEM_zero}.

\end{proof}

\end{document}
